%% =========================================================================
\section{Exercise 3: Connection to an Infinite Bus}
%% =========================================================================

Generator~A is replaced by an ideal three-phase voltage source representing an
\textbf{infinite bus} --- a bus that maintains constant terminal voltage and
constant frequency regardless of the active or reactive power exchanged with
it. This models a large national power grid to which a small generator is
connected. By definition, the infinite bus has zero impedance ($Z_\infty
\rightarrow 0$) and infinite apparent power capacity ($S_\infty \rightarrow
\infty$).

The fundamental consequence is that the infinite bus completely dominates
frequency: no action of Generator~B can alter the system frequency.
Generator~B must adapt entirely to the grid's frequency, not the other way
around.

%% ── SIMULINK MODEL DIAGRAM ───────────────────────────────────────────────
\subsection{Simulink Model}

\figref{fig:ex3_diagram} shows the Exercise~3 Simulink model. Generator~A
has been replaced by an ideal three-phase voltage source (infinite bus) that
maintains constant voltage at \pu{1.0} and constant frequency at \Hz{60}.
Generator~B connects through the synchronisation breaker at $t=\SI{5}{s}$.
The infinite bus absorbs or supplies any real and reactive power demanded by
Gen~B without any change in its terminal conditions.

\begin{figure}[H]
  \centering
  \includegraphics[width=0.75\textwidth]{EX3_figure.png}
  \caption{Exercise 3 Simulink model: Generator~B connected to an infinite bus
    (ideal voltage source) through a synchronisation breaker. The infinite bus
    maintains $V = \pu{1.0}$ and $f = \Hz{60}$ at all times.}
  \label{fig:ex3_diagram}
\end{figure}

%% ── TASK 5 ────────────────────────────────────────────────────────────────
\subsection{Task 5 -- Generator B at Lower Frequency (0.98\,pu) vs Infinite Bus}

\subsubsection{Frequency (Task 5)}

\begin{figure}[H]
  \centering
  \includegraphics[width=0.65\textwidth]{task5_frequeancy.png}
  \caption{Task 5: Frequency variation when Generator~B is paralleled with
    an infinite bus at $f_\text{set} = \pu{0.98}$. System frequency remains
    rigid at \pu{1.000}.}
  \label{fig:task5_freq}
\end{figure}

Upon breaker closure at $t=\SI{5}{s}$, the infinite bus instantaneously
forces Gen~B's frequency to \pu{1.0}. The system frequency does not change
at all --- it remains exactly \pu{1.0} throughout. This is the defining
characteristic of an infinite bus and the fundamental contrast with Task~3.

\subsubsection{Active Power (Task 5)}

\begin{figure}[H]
  \centering
  \includegraphics[width=0.65\textwidth]{task5_active_power.png}
  \caption{Task 5: Active power when Generator~B is paralleled with an
    infinite bus at $f_\text{set} = \pu{0.98}$. Gen~B absorbs power
    (severe motoring condition).}
  \label{fig:task5_active}
\end{figure}

With the grid frequency fixed at \pu{1.0} and Gen~B's setpoint at \pu{0.98},
the droop equation gives:

\begin{equation}
  P_B = \frac{f_\text{setpoint} - f_\text{grid}}{S_P}
  = \frac{0.98 - 1.0}{S_P} < 0
  \label{eq:task5_PB}
\end{equation}

Generator~B operates as a synchronous motor, absorbing power from the infinite
bus. Comparing with Task~3, where the system frequency settled at a compromise
value of $\approx\pu{1.008}$ (producing a deficit of
$1.008 - 0.98 = 0.028$\,pu), the infinite bus scenario yields a smaller
steady-state deficit of $1.000 - 0.98 = 0.020$\,pu. Although the droop-equation
deficit is numerically smaller, the practical consequence remains serious:
Gen~B cannot drag the system frequency downward as it did in Task~3, so the
motoring persists indefinitely until the operator raises the governor setpoint
or trips the breaker.

\subsubsection{Voltage and Reactive Power (Task 5)}

\begin{figure}[H]
  \centering
  \includegraphics[width=0.65\textwidth]{task5_voltage.png}
  \caption{Task 5: Terminal voltage when Generator~B is connected to an
    infinite bus at $f_\text{set} = \pu{0.98}$.}
  \label{fig:task5_voltage}
\end{figure}

\begin{figure}[H]
  \centering
  \includegraphics[width=0.65\textwidth]{task5_reactive_power.png}
  \caption{Task 5: Reactive power when Generator~B is connected to an
    infinite bus at $f_\text{set} = \pu{0.98}$.}
  \label{fig:task5_reactive}
\end{figure}

Upon breaker closure, the terminal voltage experiences a brief transient dip
caused by the circulating current between Gen~B and the infinite bus. Because
the infinite bus holds its voltage at exactly \pu{1.0}, Gen~B's AVR is the
only regulator that needs to act; the grid-side voltage is unaffected. The AVR
increases field excitation to restore $V_t$, which also raises the reactive
power output of Gen~B. In steady state, Gen~B settles at a reactive power
level determined by the difference between its internal EMF and the bus
voltage. Compared with Task~3, where both AVRs participated in voltage
recovery, the infinite bus case places the entire reactive power burden on
Gen~B's excitation system.

\subsubsection{Comparison with Task 3}

\tabref{tab:t3_vs_t5} quantifies the contrast between the two scenarios.

\begin{table}[H]
  \centering
  \small
  \caption{Task 3 vs Task 5 -- Finite Bus vs Infinite Bus ($f_\text{set,B} = \pu{0.98}$)}
  \label{tab:t3_vs_t5}
  \begin{tabular}{>{\bfseries}p{3.2cm} p{5cm} p{5cm}}
    \toprule
    \textbf{Parameter}
      & \textbf{Task 3 (Finite Bus)}
      & \textbf{Task 5 (Infinite Bus)} \\
    \midrule
    System frequency after sync
      & $\approx\pu{1.008}$ (compromise between setpoints)
      & $\pu{1.000}$ exactly (rigid, unchanged) \\[4pt]
    Gen~B active power
      & Slightly negative (weak motoring)
      & Strongly negative (severe motoring) \\[4pt]
    Freq.\ influence of Gen~B
      & Drags system freq.\ down partially
      & Zero influence on system freq. \\[4pt]
    Voltage regulation
      & Both AVRs act to restore \pu{1.0}
      & Grid holds voltage; AVR acts on Gen~B only \\[4pt]
    Engineering implication
      & Overloads Gen~A; unstable if repeated
      & Gen~B trips on overload; must raise setpoint \\
    \bottomrule
  \end{tabular}
\end{table}

%% ── TASK 6 ────────────────────────────────────────────────────────────────
\subsection{Task 6 -- Generator B at Higher Frequency (1.02\,pu) vs Infinite Bus}

\subsubsection{Frequency (Task 6)}

\begin{figure}[H]
  \centering
  \includegraphics[width=0.65\textwidth]{task6_frequeancy.png}
  \caption{Task 6: Frequency variation when Generator~B is paralleled with
    an infinite bus at $f_\text{set} = \pu{1.02}$. System frequency remains
    rigid at \pu{1.000}.}
  \label{fig:task6_freq}
\end{figure}

Generator~B is spun up to \pu{1.02} before synchronisation. Upon breaker
closure, the infinite bus immediately locks Gen~B to \pu{1.0}. The frequency
does not rise --- it remains exactly \pu{1.0}. This is the key contrast with
Task~4, where Gen~A (a finite machine) was pulled upward in frequency by Gen~B.

\subsubsection{Active Power (Task 6)}

\begin{figure}[H]
  \centering
  \includegraphics[width=0.65\textwidth]{task6_active_power.png}
  \caption{Task 6: Active power when Generator~B is paralleled with an
    infinite bus at $f_\text{set} = \pu{1.02}$. Gen~B injects power into
    the grid.}
  \label{fig:task6_active}
\end{figure}

With Gen~B's setpoint above the infinite bus frequency:

\begin{equation}
  P_B = \frac{f_\text{setpoint} - f_\text{grid}}{S_P}
  = \frac{1.02 - 1.0}{S_P} > 0
  \label{eq:task6_PB}
\end{equation}

Generator~B injects positive active power into the infinite bus. The grid
absorbs this power effortlessly without any frequency perturbation.

\subsubsection{Voltage and Reactive Power (Task 6)}

\begin{figure}[H]
  \centering
  \includegraphics[width=0.65\textwidth]{task6_voltage.png}
  \caption{Task 6: Terminal voltage when Generator~B is connected to an
    infinite bus at $f_\text{set} = \pu{1.02}$.}
  \label{fig:task6_voltage}
\end{figure}

\begin{figure}[H]
  \centering
  \includegraphics[width=0.65\textwidth]{task6_reactive_power.png}
  \caption{Task 6: Reactive power when Generator~B is connected to an
    infinite bus at $f_\text{set} = \pu{1.02}$.}
  \label{fig:task6_reactive}
\end{figure}

The voltage transient at breaker closure in Task~6 is a brief overshoot rather
than a dip, because Gen~B's rotor was leading the infinite bus in phase angle
at the moment of connection. The AVR damps this overshoot within one to two
seconds. In steady state, Gen~B exports reactive power set by the difference
between its internal EMF (driven by the AVR setpoint) and the fixed bus voltage
of \pu{1.0}. The infinite bus absorbs any surplus reactive power without any
voltage deviation, which contrasts with Task~4 where both generators shared the
reactive power adjustment. This clean voltage recovery confirms that the
AVR--exciter loop operates independently of the P--f loop, consistent with the
decoupled control design discussed in Exercise~1.

\subsubsection{Comparison with Task 4}

\tabref{tab:t4_vs_t6} quantifies the contrast.

\begin{table}[H]
  \centering
  \small
  \caption{Task 4 vs Task 6 -- Finite Bus vs Infinite Bus ($f_\text{set,B}: \pu{1.02}$)}
  \label{tab:t4_vs_t6}
  \begin{tabular}{>{\bfseries}p{3.2cm} p{5cm} p{5cm}}
    \toprule
    \textbf{Parameter}
      & \textbf{Task 4 (Finite Bus)}
      & \textbf{Task 6 (Infinite Bus)} \\
    \midrule
    System frequency after sync
      & Raised by Gen~B (droop-weighted average)
      & $\pu{1.000}$ exactly (cannot be raised) \\[4pt]
    Gen~A / Grid response
      & Gen~A reduces output significantly
      & Grid absorbs all Gen~B output silently \\[4pt]
    Gen~B active power
      & Majority of load (generation)
      & Fixed by setpoint offset from grid \\[4pt]
    Voltage behaviour
      & Spike then \pu{1.0}
      & Spike then \pu{1.0} (identical mechanism) \\[4pt]
    Freq.\ influence of Gen~B
      & Raises system frequency
      & Cannot raise grid frequency \\[4pt]
    Engineering implication
      & Load transferred; useful commissioning
      & Normal grid-tie generation; power sold \\
    \bottomrule
  \end{tabular}
\end{table}

%% ── GRAND COMPARISON ─────────────────────────────────────────────────────
\subsection{Grand Comparison -- All Tasks}

The following comparison across all four parallel-operation tasks (Tasks~3--6)
shows how the bus type (finite vs infinite) and the governor setpoint
(under-frequency vs over-frequency) interact to determine system behaviour.

Key observations are summarised in \tabref{tab:grand_comparison}.

\begin{table}[H]
  \centering
  \small
  \caption{Grand Comparison of Tasks 3--6}
  \label{tab:grand_comparison}
  \begin{tabular}{>{\bfseries}p{2.4cm} p{2.4cm} p{2.4cm} p{2.4cm} p{2.4cm}}
    \toprule
    \textbf{Parameter}
      & \textbf{Task 3} (Finite, 0.98)
      & \textbf{Task 4} (Finite, 1.02)
      & \textbf{Task 5} (Infinite, 0.98)
      & \textbf{Task 6} (Infinite, 1.02) \\
    \midrule
    System freq.\ after sync
      & $\approx\pu{1.008}$ (compromise)
      & $\approx\pu{1.015}$ (compromise)
      & \pu{1.000} (rigid)
      & \pu{1.000} (rigid) \\[4pt]
    Gen~B active power
      & Negative (motoring)
      & Positive (generation)
      & Negative (motoring)
      & Positive (generation) \\[4pt]
    Gen~B freq.\ influence
      & Pulls system freq.\ down
      & Pulls system freq.\ up
      & None
      & None \\[4pt]
    Voltage regulation
      & Both AVRs act
      & Both AVRs act
      & Grid holds $V$; Gen~B AVR only
      & Grid holds $V$; Gen~B AVR only \\[4pt]
    Engineering implication
      & Overloads Gen~A
      & Load transferred successfully
      & Gen~B trips on overload
      & Normal grid-tie operation \\
    \bottomrule
  \end{tabular}
\end{table}

The governor setpoint relative to the system frequency determines whether a
generator motors or generates in every scenario. The bus type (finite vs
infinite) controls whether the oncoming machine can influence system frequency:
a finite bus allows a droop-weighted compromise, while an infinite bus holds
frequency immovable at \pu{1.0}.
